% !TeX program = xelatex
% HomeworkReviewOS | Ярослав Гаврилов | 22.10.2026
% Компиляция в Overleaf: XeLaTeX, формат A4.
\documentclass[12pt,a4paper]{article}
\usepackage[a4paper,left=22mm,right=22mm,top=22mm,bottom=23mm]{geometry}
\usepackage{fontspec}
\setmainfont{DejaVu Serif}
\usepackage{polyglossia}
\setdefaultlanguage{russian}
\usepackage{amsmath,amssymb,mathtools}
\usepackage{microtype}
\usepackage{enumitem}
\usepackage{fancyhdr}
\usepackage{titlesec}
\usepackage[most]{tcolorbox}
\usepackage{xcolor}
\definecolor{ink}{HTML}{163047}
\definecolor{muted}{HTML}{526577}
\definecolor{accent}{HTML}{0F6C85}
\definecolor{soft}{HTML}{F2F8FA}
\definecolor{line}{HTML}{D7E5EC}
\color{ink}
\titleformat{\section}{\large\bfseries\color{ink}}{}{0pt}{}
\titlespacing*{\section}{0pt}{18pt}{7pt}
\setlength{\parindent}{0pt}
\setlength{\parskip}{7pt}
\setlength{\headheight}{15pt}
\pagestyle{fancy}
\fancyhf{}
\fancyhead[L]{\small\color{muted}Разбор домашней работы}
\fancyhead[R]{\small\color{muted}22 октября 2026 г.}
\fancyfoot[C]{\small\color{muted}\thepage}
\renewcommand{\headrulewidth}{0.3pt}
\renewcommand{\footrulewidth}{0pt}
\tcbset{enhanced,breakable,boxrule=.45pt,arc=1.4mm,left=3mm,right=3mm,top=2.5mm,bottom=2.5mm,colframe=line,colback=soft}
\newtcolorbox{practice}{title={\textbf{Задача на закрепление}},fonttitle=\normalsize\bfseries,coltitle=ink,colbacktitle=soft,attach boxed title to top left={xshift=2mm,yshift=-1mm},boxed title style={frame hidden,interior hidden},before skip=11pt,after skip=12pt}
\newcommand{\Step}[1]{\par\textbf{#1}\quad}
\begin{document}
{\Large\bfseries\color{ink}Проверка домашней работы}\par
{\large Ярослав Гаврилов}\par
{\small\color{muted}Тема: тригонометрия, свойства степеней и показательные уравнения}\par
\vspace{2mm}
\begin{tcolorbox}
\textbf{На что обратить внимание.} В работе обнаружены две ошибки при преобразовании степеней (задания 6 и 7). В задании 10 вычисления верны, однако окончательный ответ оформлен в другом виде, чем требуется в условии. Ниже приведены разборы только этих моментов.
\end{tcolorbox}

\section*{Задание 6. Не потерять внешнюю степень}
\textbf{Исходное выражение:}
\[
\left(\frac{3^{1/2}\cdot 3^{1/3}}{\sqrt[6]{3}}\right)^3.
\]
\Step{Где возникла ошибка.} В твоих вычислениях степень внутри скобок упрощена до \(3^{2/3}\), затем записан ответ \(\sqrt[3]{9}\). В этом переходе пропущена внешняя третья степень, которая относится ко \emph{всей дроби}. Полученный результат соответствует выражению \emph{внутри} скобок, но ещё не всему заданию.

\Step{Как решить правильно.} Сначала представим корень как дробную степень:
\[
\sqrt[6]{3}=3^{1/6}.
\]
При умножении степеней с одинаковыми основаниями показатели складываются. При делении показатели вычитаются. Внутри скобок получаем
\[
\frac{3^{1/2}\cdot3^{1/3}}{3^{1/6}}
=3^{\,1/2+1/3-1/6}
=3^{\,3/6+2/6-1/6}
=3^{2/3}.
\]
Теперь обязательно учитываем третью степень, стоящую \emph{снаружи}:
\[
\left(3^{2/3}\right)^3
=3^{\,(2/3)\cdot3}
=3^2=\boxed{9}.
\]
\textbf{Вывод.} При возведении степени в степень показатели \emph{перемножаются}. Полезно сначала упростить внутреннее выражение, затем отдельно обработать внешнюю степень.
\begin{practice}
Вычисли, последовательно применяя свойства степеней:
\[
\left(\frac{5^{1/2}\cdot5^{1/4}}{\sqrt[4]{5}}\right)^4.
\]
\end{practice}

\section*{Задание 7. Раскрытие скобок в показателях}
\textbf{Исходное выражение:}
\[
\frac{(0{,}25)^{\sqrt5-2}}{2^{\,2-\sqrt5}}.
\]
\Step{Где возникла ошибка.} В записи преобразований изменились показатели степеней. При замене \(0{,}25\) степенью двойки потерян результат умножения числа \(-2\) на свободный член \(-2\); в показателе знаменателя также записано другое выражение. Из-за этого в конце получилось \(16\). Здесь каждый знак существенно влияет на результат.

\Step{Как решить правильно.} Десятичную дробь приводим к степени с основанием \(2\):
\[
0{,}25=\frac14=2^{-2}.
\]
Возводим степень в степень. По распределительному свойству умножения:
\[
(0{,}25)^{\sqrt5-2}
=(2^{-2})^{\sqrt5-2}
=2^{-2(\sqrt5-2)}
=2^{\,4-2\sqrt5}.
\]
Обрати внимание: \(-2\cdot(-2)=+4\). В знаменателе сохраняем исходный показатель \(2-\sqrt5\). При делении степеней с основанием \(2\) вычитаем \emph{весь} показатель знаменателя:
\begin{align*}
\frac{2^{\,4-2\sqrt5}}{2^{\,2-\sqrt5}}
&=2^{\,(4-2\sqrt5)-(2-\sqrt5)}\\
&=2^{\,4-2\sqrt5-2+\sqrt5}\\
&=\boxed{2^{\,2-\sqrt5}}.
\end{align*}
\Step{Как проверить себя.} Поскольку \(\sqrt5>2\), показатель \(2-\sqrt5\) отрицателен. Следовательно, найденная степень двойки меньше единицы. Ответ \(16\) этому свойству противоречит: такая быстрая проверка помогает обнаружить ошибку ещё до переписывания решения.
\begin{practice}
Упрости выражение до одной степени с основанием \(2\):
\[
\frac{(0{,}25)^{\sqrt3-1}}{2^{\,1-\sqrt3}}.
\]
\end{practice}

\clearpage
\section*{Задание 10. Соблюдать требуемый формат ответа}
\textbf{Условие:} \(Pv^{3/2}=C\), где \(P=2{,}5\cdot10^3\), \(C=20\), \(v>0\). Ответ требуется записать \emph{десятичной дробью}.

\Step{Что произошло.} Ты правильно получил \(v=\dfrac1{25}\). Значение найдено верно, но финальная запись дана обыкновенной дробью. В задании специально оговорён десятичный вид ответа, поэтому решение нужно довести до этого формата.

\Step{Полная запись.} Подставляем известные значения и делим обе части на \(2500\):
\[
2500v^{3/2}=20,
\qquad v^{3/2}=\frac{20}{2500}=\frac1{125}=5^{-3}.
\]
Чтобы найти \(v\), возводим обе части в степень \(\frac23\). Это допустимо, поскольку \(v>0\):
\[
v=\left(5^{-3}\right)^{2/3}=5^{-2}=\frac1{25}.
\]
Переводим обыкновенную дробь в десятичную:
\[
\frac1{25}=\frac4{100}=\boxed{0{,}04}.
\]
\textbf{Вывод.} В конце решения ещё раз прочитай требование к форме ответа: число может быть вычислено верно, но ответ нужно записать в указанном виде.
\begin{practice}
Величины \(P\) и \(v>0\) связаны формулой \(Pv^{3/2}=C\). Найди \(v\), если \(P=400\), \(C=50\). Запиши ответ десятичной дробью.
\end{practice}

\vfill
\begin{center}
\footnotesize\color{muted}Лёвин Артём Александрович эксклюзивно для Ярослава Гаврилова
\end{center}
\end{document}
